\section{AHSME 1972}
        \begin{enumerate}
        	\item (\textit{AHSME 1972}) The lengths in inches of the three sides of each of four triangles $I, II, III$, and $IV$ are as follows: 


 $\begin{array}{rlrl} \hbox{I}& 3,\ 4,\ \hbox{and}\ 5\qquad & \hbox{III}& 7,\ 24,\ \hbox{and}\ 25\\  \hbox{II}& 4,\ 7\frac{1}{2},\ \hbox{and}\ 8\frac{1}{2}\qquad & \hbox{IV}& 3\frac{1}{2},\ 4\frac{1}{2},\ \hbox{and}\ 5\frac{1}{2}.\end{array}$


 Of these four given triangles, the only right triangles are 


 $\text{(A)  I and II}\qquad  \text{(B)  I and III}\qquad \text{(C)  I and IV} \quad \\ \text{(D)  I, II, and III} \qquad \text{(E)  I, II, and IV }$


	\item (\textit{AHSME 1972}) If a dealer could get his goods for $8$\% less while keeping his selling price fixed, 
his profit, based on cost, would be increased to $(x+10)$\% from his present profit of $x$\%, which is


 $\textbf{(A) }12\%\qquad \textbf{(B) }15\%\qquad \textbf{(C) }30\%\qquad \textbf{(D) }50\%\qquad  \textbf{(E) }75\%$


	\item (\textit{AHSME 1972}) If $x=\dfrac{1-i\sqrt{3}}{2}$ where $i=\sqrt{-1}$, then $\dfrac{1}{x^2-x}$ is equal to


 $\textbf{(A) }-2\qquad \textbf{(B) }-1\qquad \textbf{(C) }1+i\sqrt{3}\qquad \textbf{(D) }1\qquad  \textbf{(E) }2$


	\item (\textit{AHSME 1972}) The number of solutions to $\{1,~2\}\subseteq~X~\subseteq~\{1,~2,~3,~4,~5\}$, where $X$ is a subset of $\{1,~2,~3,~4,~5\}$ is


 $\textbf{(A) }2\qquad \textbf{(B) }4\qquad \textbf{(C) }6\qquad \textbf{(D) }8\qquad  \textbf{(E) }\text{None of these}$


	\item (\textit{AHSME 1972}) From among $2^{1/2}, 3^{1/3}, 8^{1/8}, 9^{1/9}$ those which have the greatest and the next to the greatest values, in that order, are 


 $\textbf{(A) } 3^{1/3},\ 2^{1/2}\quad  \textbf{(B) } 3^{1/3},\ 8^{1/8}\quad  \textbf{(C) } 3^{1/3},\ 9^{1/9}\quad  \textbf{(D) } 8^{1/8},\ 9^{1/9}\quad \\  \text{(E) None of these}$


	\item (\textit{AHSME 1972}) If $3^{2x}+9=10\left(3^{x}\right)$, then the value of $(x^2+1)$ is 


 $\textbf{(A) }1\text{ only}\qquad \textbf{(B) }5\text{ only}\qquad \textbf{(C) }1\text{ or }5\qquad \textbf{(D) }2\qquad \textbf{(E) }10$


	\item (\textit{AHSME 1972}) If $yz:zx:xy=1:2:3$, then $\dfrac{x}{yz}:\dfrac{y}{zx}$ is equal to


 $\textbf{(A) }3:2\qquad \textbf{(B) }1:2\qquad \textbf{(C) }1:4\qquad \textbf{(D) }2:1\qquad  \textbf{(E) }4:1$


	\item (\textit{AHSME 1972}) If $|x-\log y|=x+\log y$ where $x$ and $\log y$ are real, then


 $\textbf{(A) }x=0\qquad \textbf{(B) }y=1\qquad \textbf{(C) }x=0\text{ and }y=1\qquad\\ \textbf{(D) }x(y-1)=0\qquad  \textbf{(E) }\text{None of these}$


	\item (\textit{AHSME 1972}) Ann and Sue bought identical boxes of stationery. Ann used hers to write $1$-sheet letters and Sue used hers to write $3$-sheet letters. 
Ann used all the envelopes and had $50$ sheets of paper left, while Sue used all of the sheets of paper and had $50$ envelopes left. 
The number of sheets of paper in each box was 


 $\textbf{(A) } 150 \qquad \textbf{(B) } 125 \qquad \textbf{(C) } 120 \qquad \textbf{(D) } 100 \qquad     \textbf{(E) } 80 \qquad$


	\item (\textit{AHSME 1972}) For $x$ real, the inequality $1\le |x-2|\le 7$ is equivalent to


 $\textbf{(A) }x\le 1\text{ or }x\ge 3\qquad \textbf{(B) }1\le x\le 3\qquad \textbf{(C) }-5\le x\le 9\qquad \\ \textbf{(D) }-5\le x\le 1\text{ or }3\le x\le 9\qquad  \textbf{(E) }-6\le x\le 1\text{ or }3\le x\le 10$


	\item (\textit{AHSME 1972}) The value(s) of $y$ for which the following pair of equations $x^2+y^2-16=0\text{ and }x^2-3y+12=0$ may have a real common solution, are


 $\textbf{(A) }4\text{ only}\qquad \textbf{(B) }-7,~4\qquad \textbf{(C) }0,~4\qquad \textbf{(D) }\text{no }y\qquad  \textbf{(E) }\text{all }y$


	\item (\textit{AHSME 1972}) The number of cubic feet in the volume of a cube is the same as the number of square inches in its surface area. The length of the edge expressed as a number of feet is


 $\textbf{(A) }6\qquad \textbf{(B) }864\qquad \textbf{(C) }1728\qquad \textbf{(D) }6\times 1728\qquad  \textbf{(E) }2304$


	\item (\textit{AHSME 1972}) 
                    \begin{figure}[H]
                        \centering
                        \includegraphics[width=0.5\textwidth]{Problem/AHSME_I_1972_13_Problem_Figure_1.png}
                    \end{figure}
                    


 Inside square $ABCD$ (See figure) with sides of length $12$ inches, segment $AE$ is drawn where $E$ is the point on $DC$ which is $5$ inches from $D$. 
The perpendicular bisector of $AE$ is drawn and intersects $AE, AD$, and $BC$ at points $M, P$, and $Q$ respectively. The ratio of segment $PM$ to $MQ$ is


 $\textbf{(A) }5:12\qquad \textbf{(B) }5:13\qquad \textbf{(C) }5:19\qquad \textbf{(D) }1:4\qquad  \textbf{(E) }5:21$


	\item (\textit{AHSME 1972}) A triangle has angles of $30^\circ$ and $45^\circ$. If the side opposite the $45^\circ$ angle has length $8$, then the side opposite the $30^\circ$ angle has length


 $\textbf{(A) }4\qquad \textbf{(B) }4\sqrt{2}\qquad \textbf{(C) }4\sqrt{3}\qquad \textbf{(D) }4\sqrt{6}\qquad  \textbf{(E) }6$


	\item (\textit{AHSME 1972}) A contractor estimated that one of his two bricklayers would take $9$ hours to build a certain wall and the other $10$ hours. 
However, he knew from experience that when they worked together, their combined output fell by $10$ bricks per hour. 
Being in a hurry, he put both men on the job and found that it took exactly 5 hours to build the wall. The number of bricks in the wall was


 $\textbf{(A) }500\qquad \textbf{(B) }550\qquad \textbf{(C) }900\qquad \textbf{(D) }950\qquad  \textbf{(E) }960$


	\item (\textit{AHSME 1972}) There are two positive numbers that may be inserted between $3$ and $9$ such that the first three are in geometric progression 
while the last three are in arithmetic progression. The sum of those two positive numbers is


 $\textbf{(A) }13\textstyle\frac{1}{2}\qquad \textbf{(B) }11\frac{1}{4}\qquad \textbf{(C) }10\frac{1}{2}\qquad \textbf{(D) }10\qquad \textbf{(E) }9\frac{1}{2}$


	\item (\textit{AHSME 1972}) A piece of string is cut in two at a point selected at random. The probability that the longer piece is at least x times as large as the shorter piece is


 $\textbf{(A) }\frac{1}{2}\qquad \textbf{(B) }\frac{2}{x}\qquad \textbf{(C) }\frac{1}{x+1}\qquad \textbf{(D) }\frac{1}{x}\qquad \textbf{(E) }\frac{2}{x+1}$


	\item (\textit{AHSME 1972}) Let $ABCD$ be a trapezoid with the measure of base $AB$ twice that of base $DC$, and let $E$ be the point of intersection of the diagonals. 
If the measure of diagonal $AC$ is $11$, then that of segment $EC$ is equal to


 $\textbf{(A) }3\textstyle\frac{2}{3}\qquad \textbf{(B) }3\frac{3}{4}\qquad \textbf{(C) }4\qquad \textbf{(D) }3\frac{1}{2}\qquad \textbf{(E) }3$


	\item (\textit{AHSME 1972}) The sum of the first $n$ terms of the sequence $1,~(1+2),~(1+2+2^2),~\dots ~(1+2+2^2+\dots +2^{n-1})$ in terms of $n$ is


 $\textbf{(A) }2^n\qquad \textbf{(B) }2^n-n\qquad \textbf{(C) }2^{n+1}-n\qquad \textbf{(D) }2^{n+1}-n-2\qquad  \textbf{(E) }n\cdot 2^n$


	\item (\textit{AHSME 1972}) If $\tan x=\dfrac{2ab}{a^2-b^2}$ where $a>b>0$ and $0^\circ <x<90^\circ$, then $\sin x$ is equal to


 $\textbf{(A) }\frac{a}{b}\qquad \textbf{(B) }\frac{b}{a}\qquad \textbf{(C) }\frac{\sqrt{a^2-b^2}}{2a}\qquad \textbf{(D) }\frac{\sqrt{a^2-b^2}}{2ab}\qquad \textbf{(E) }\frac{2ab}{a^2+b^2}$


	\item (\textit{AHSME 1972}) 
                    \begin{figure}[H]
                        \centering
                        \includegraphics[width=0.5\textwidth]{Problem/AHSME_I_1972_21_Problem_Figure_1.png}
                    \end{figure}
                    


 If the sum of the measures in degrees of angles $A,~B,~C,~D,~E$ and $F$ in the figure above is $90n$, then $n$ is equal to


 $\textbf{(A) }2\qquad \textbf{(B) }3\qquad \textbf{(C) }4\qquad \textbf{(D) }5\qquad \textbf{(E) }6$


	\item (\textit{AHSME 1972}) If $a\pm bi~(b\neq 0)$ are imaginary roots of the equation $x^3+qx+r=0$ where $a,~b,~q$, and $r$ are real numbers, then $q$ in terms of $a$ and $b$ is


 $\textbf{(A) }a^2+b^2\qquad \textbf{(B) }2a^2-b^2\qquad \textbf{(C) }b^2-a^2\qquad \textbf{(D) }b^2-2a^2\qquad  \textbf{(E) }b^2-3a^2$


	\item (\textit{AHSME 1972}) 
                    \begin{figure}[H]
                        \centering
                        \includegraphics[width=0.5\textwidth]{Problem/AHSME_I_1972_23_Problem_Figure_1.png}
                    \end{figure}
                    


 The radius of the smallest circle containing the symmetric figure composed of the 3 unit squares shown above is


 $\textbf{(A) }\sqrt{2}\qquad \textbf{(B) }\sqrt{1.25}\qquad \textbf{(C) }1.25\qquad \textbf{(D) }\frac{5\sqrt{17}}{16}\qquad  \textbf{(E) }\text{None of these}$


	\item (\textit{AHSME 1972}) A man walked a certain distance at a constant rate. If he had gone $\textstyle\frac{1}{2}$ mile per hour faster, he would have walked the distance in four-fifths of the time; if he had gone $\textstyle\frac{1}{2}$ mile per hour slower, he would have been $2\textstyle\frac{1}{2}$ hours longer on the road. The distance in miles he walked was


 $\textbf{(A) }13\textstyle\frac{1}{2}\qquad \textbf{(B) }15\qquad \textbf{(C) }17\frac{1}{2}\qquad \textbf{(D) }20\qquad  \textbf{(E) }25$


	\item (\textit{AHSME 1972}) Inscribed in a circle is a quadrilateral having sides of lengths $25,~39,~52$, and $60$ taken consecutively. The diameter of this circle has length


 $\textbf{(A) }62\qquad \textbf{(B) }63\qquad \textbf{(C) }65\qquad \textbf{(D) }66\qquad  \textbf{(E) }69$


	\item (\textit{AHSME 1972}) 
                    \begin{figure}[H]
                        \centering
                        \includegraphics[width=0.5\textwidth]{Problem/AHSME_I_1972_26_Problem_Figure_1.png}
                    \end{figure}
                    


 In the circle above, $M$ is the midpoint of arc $CAB$ and segment $MP$ is perpendicular to chord $AB$ at $P$. 
If the measure of chord $AC$ is $x$ and that of segment $AP$ is $(x+1)$, then segment $PB$ has measure equal to


 $\textbf{(A) }3x+2\qquad \textbf{(B) }3x+1\qquad \textbf{(C) }2x+3\qquad \textbf{(D) }2x+2\qquad  \textbf{(E) }2x+1$


	\item (\textit{AHSME 1972}) If the area of $\triangle ABC$ is $64$ square units and the geometric mean (mean proportional) 
between sides $AB$ and $AC$ is $12$ inches, then $\sin A$ is equal to


 $\textbf{(A) }\dfrac{\sqrt{3}}{2}\qquad \textbf{(B) }\frac{3}{5}\qquad \textbf{(C) }\frac{4}{5}\qquad \textbf{(D) }\frac{8}{9}\qquad \textbf{(E) }\frac{15}{17}$


	\item (\textit{AHSME 1972}) A circular disc with diameter $D$ is placed on an $8\times 8$ checkerboard with width $D$ so that the centers coincide. The number of checkerboard squares which are completely covered by the disc is


 $\textbf{(A) }48\qquad \textbf{(B) }44\qquad \textbf{(C) }40\qquad \textbf{(D) }36\qquad  \textbf{(E) }32$


	\item (\textit{AHSME 1972}) If $f(x)=\log \left(\frac{1+x}{1-x}\right)$ for $-1<x<1$, then $f\left(\frac{3x+x^3}{1+3x^2}\right)$ in terms of $f(x)$ is


 $\textbf{(A) }-f(x)\qquad \textbf{(B) }2f(x)\qquad \textbf{(C) }3f(x)\qquad \textbf{(D) }\left[f(x)\right]^2\qquad \\ \textbf{(E) }[f(x)]^3-f(x)$


	\item (\textit{AHSME 1972}) 
                    \begin{figure}[H]
                        \centering
                        \includegraphics[width=0.5\textwidth]{Problem/AHSME_I_1972_30_Problem_Figure_1.png}
                    \end{figure}
                    


 A rectangular piece of paper 6 inches wide is folded as in the diagram so that one corner touches the opposite side. The length in inches of the crease L in terms of angle $\theta$ is


 $\textbf{(A) }3\sec ^2\theta\csc\theta\qquad \textbf{(B) }6\sin\theta\sec\theta\qquad \textbf{(C) }3\sec\theta\csc\theta\qquad \textbf{(D) }6\sec\theta\csc^2\theta\qquad \textbf{(E) }\text{None of these}$


	\item (\textit{AHSME 1972}) When the number $2^{1000}$ is divided by $13$, the remainder in the division is


 $\textbf{(A) }1\qquad \textbf{(B) }2\qquad \textbf{(C) }3\qquad \textbf{(D) }7\qquad  \textbf{(E) }11$


	\item (\textit{AHSME 1972}) 
                    \begin{figure}[H]
                        \centering
                        \includegraphics[width=0.5\textwidth]{Problem/AHSME_I_1972_32_Problem_Figure_1.png}
                    \end{figure}
                    


 Chords $AB$ and $CD$ in the circle above intersect at E and are perpendicular to each other. 
If segments $AE, EB$, and $ED$ have measures $2, 3$, and $6$ respectively, then the length of the diameter of the circle is


 $\textbf{(A) }4\sqrt{5}\qquad \textbf{(B) }\sqrt{65}\qquad \textbf{(C) }2\sqrt{17}\qquad \textbf{(D) }3\sqrt{7}\qquad  \textbf{(E) }6\sqrt{2}$


	\item (\textit{AHSME 1972}) The minimum value of the quotient of a (base ten) number of three different non-zero digits divided by the sum of its digits is


 $\textbf{(A) }9.7\qquad \textbf{(B) }10.1\qquad \textbf{(C) }10.5\qquad \textbf{(D) }10.9\qquad  \textbf{(E) }20.5$


	\item (\textit{AHSME 1972}) Three times Dick's age plus Tom's age equals twice Harry's age. 
Double the cube of Harry's age is equal to three times the cube of Dick's age added to the cube of Tom's age. 
Their respective ages are relatively prime to each other. The sum of the squares of their ages is


 $\textbf{(A) }42\qquad \textbf{(B) }46\qquad \textbf{(C) }122\qquad \textbf{(D) }290\qquad  \textbf{(E) }326$


	\item (\textit{AHSME 1972}) 
                    \begin{figure}[H]
                        \centering
                        \includegraphics[width=0.5\textwidth]{Problem/AHSME_I_1972_35_Problem_Figure_1.png}
                    \end{figure}
                    


 Equilateral triangle $ABP$ (see figure) with side $AB$ of length $2$ inches is placed inside square $AXYZ$ with side of length $4$ 
inches so that $B$ is on side $AX$. The triangle is rotated clockwise about $B$, then $P$, and so on along the sides of the square 
until $P$ returns to its original position. The length of the path in inches traversed by vertex $P$ is equal to


 $\textbf{(A) }20\pi/3\qquad \textbf{(B) }32\pi/3\qquad \textbf{(C) }12\pi\qquad \textbf{(D) }40\pi/3\qquad  \textbf{(E) }15\pi$


\end{enumerate}